\documentclass[11pt]{article}
\usepackage[utf8]{inputenc}
\usepackage{tgpagella} % Palatino-clone as main font
\usepackage[a4paper,margin=3cm]{geometry}
\author{Bernd Kiefer}
\date{\today}
\title{}
%\tableofcontents
\pagestyle{empty}
\begin{document}
{\LARGE\bf{}Map Modeling Guidelines}\\
\begin{enumerate}
\item set \texttt{level} attribute for \textbf{every} node (obligatory! test in map loading code!)
\item According to Simple Indoor Tagging
\begin{itemize}
\item \texttt{indoor=room/area/corridor} marks an \textbf{area}
\item \texttt{indoor=wall} marks a \textbf{way}, in most places part of the above
\item stairs/elevators are \texttt{indoor=room} \textbf{and} \texttt{stairs=yes} resp.~\texttt{highway=elevator}
\end{itemize}
\item Doors: \texttt{node} with \texttt{door=yes}
\item For navigation purposes:
\begin{itemize}
\item way(s) with \texttt{indoor=highway} and \texttt{highway=(corridor|stair|elevator|ramp)}
\item connect those highways to the doors nearby
\item when adding this for a stair or an elevator, you may have nodes at the
  same location, but with different \texttt{level} values if the staircase
  stretches over more than two levels.
\item for large areas, such as lobbies, add multiple highways to avoid to guide
  the person to middle of the room, and then back
\end{itemize}
\item Add a relation, containing all nodes and ways, with

  \texttt{
  \begin{tabular}{|c@{ = }c|}\hline
    type&building\\\hline
    name&$\langle name\rangle$\\\hline
    latitude&$\langle lat \rangle$\\\hline
    longitude&$\langle lon \rangle$\\\hline
  \end{tabular}}
\end{enumerate}
\end{document}